\chapter{The Rauch Comparison Theorem}
\label{chap:dc-rauch-comparison-theorem}

\section{The Theorem of Rauch}

For a normalized geodesic $\gamma$ and a Jacobi field $J$ with
$J(0)=0$, $|J'(0)|=1$, and $J'(0)\perp\gamma'(0)$, the local expansion is
\[
 |J(t)|=t-\frac{K}{6}t^3+R,\qquad \lim_{t\to0}\frac{R}{t^3}=0,
\]
where $K$ is the sectional curvature of the plane spanned by $\gamma'(0)$ and
$J'(0)$. Thus larger curvature produces smaller transverse Jacobi fields near
$t=0$. Rauch's theorem gives a global comparison: if $\tilde\gamma$ has no
conjugate points and its sectional curvature in planes containing
$\tilde\gamma'$ is at least that of the corresponding planes along $\gamma$,
then normalized Jacobi fields with matching initial data satisfy
$|\tilde J(t)|\leq |J(t)|$. In dimension two this reduces to Sturm comparison for
\[
 f''+Kf=0,\qquad \tilde f''+\widetilde K\tilde f=0,
\]
with $f(0)=\tilde f(0)=0$ and equal positive initial derivatives.

\begin{lemma}[Factorization of a function vanishing at the origin and index form]
  \label{lem:dc-ch10-2-1}
  \dcref{ch10:2.1}
  If $h:[0,1]\to\mathbb R$ is differentiable and $h(0)=0$, there is a differentiable
  $\phi:[0,1]\to\mathbb R$ with $\phi(0)=h'(0)$ and $h(t)=t\phi(t)$. For a geodesic
  $\gamma:[0,a]\to M$ and a piecewise differentiable field $V$ along it, define
  \[
  I_{t_0}(V,V)=\int_0^{t_0}\bigl(\langle V',V'\rangle-
  \langle R(\gamma',V)\gamma',V\rangle\bigr)\,dt.
  \]
\end{lemma}

\begin{lemma}[Index comparison for fields vanishing at both endpoints]
  \label{lem:dc-ch10-2-2}
  \uses{lem:dc-ch10-2-1}
  \dcref{ch10:2.2}
  Let $\gamma:[0,a]\to M$ have no conjugate point to $\gamma(0)$ on $(0,a]$.
  If $J$ is a normal Jacobi field, $V$ a piecewise differentiable normal field,
  $J(0)=V(0)=0$, and $J(t_0)=V(t_0)$ for $t_0\in(0,a]$, then
  \[
  I_{t_0}(J,J)\leq I_{t_0}(V,V),
  \]
  with equality exactly when $V=J$ on $[0,t_0]$.
\end{lemma}
\begin{proof}
  \sketch
  Choose a basis $J_1,\ldots,J_{n-1}$ of normal Jacobi fields vanishing at $0$.
  Absence of conjugate points makes $J_i(t)$ a basis of $\gamma'(t)^\perp$ for
  $t>0$. Write $V=\sum_i f_iJ_i$; the factorization lemma extends the coefficients
  piecewise differentiably to $0$. The Wronskian identity
  \[
  \langle J_i',J_j\rangle=\langle J_i,J_j'\rangle
  \]
  follows because its derivative is zero and its value at $0$ is zero. Integration
  by parts therefore gives
  \[
  I_{t_0}(V,V)=\Big\langle\sum_i f_iJ_i,\sum_j f_jJ_j'\Big\rangle(t_0)
  +\int_0^{t_0}\Big|\sum_i f_i'J_i\Big|^2dt,
  \]
  while $I_{t_0}(J,J)$ is the same boundary term with constant coefficients
  $\alpha_i$. Since $V(t_0)=J(t_0)$, $f_i(t_0)=\alpha_i$, and hence
  \[
  I_{t_0}(V,V)=I_{t_0}(J,J)+\int_0^{t_0}\Big|\sum_i f_i'J_i\Big|^2dt.
  \]
  The inequality and its equality criterion follow from linear independence for
  $t>0$.
\end{proof}

\begin{theorem}[Rauch comparison]
  \label{thm:dc-ch10-2-3}
  \uses{lem:dc-ch10-2-2, prop:dc-ch5-3-6}
  \dcref{ch10:2.3}
  Let $\gamma:[0,a]\to M^n$ and $\tilde\gamma:[0,a]\to\widetilde M^{n+k}$ be
  geodesics with $|\gamma'|=|\tilde\gamma'|$. Let $J,\tilde J$ be Jacobi fields
  satisfying
  \[
  J(0)=\tilde J(0)=0,\quad
  \langle J'(0),\gamma'(0)\rangle=\langle\tilde J'(0),\tilde\gamma'(0)\rangle,
  \quad |J'(0)|=|\tilde J'(0)|.
  \]
  Assume that $\tilde\gamma$ has no conjugate points on $(0,a]$ and that
  \[
  \widetilde K(\tilde x,\tilde\gamma'(t))\geq K(x,\gamma'(t))
  \]
  for every $t$ and all tangent vectors $x,\tilde x$ (with the sectional
  curvatures taken in the planes they span with the geodesic velocities). Then
  $|\tilde J(t)|\leq |J(t)|$ on $[0,a]$. If equality holds at some
  $t_0\in(0,a]$, then
  \[
  \widetilde K(\tilde J(t),\tilde\gamma'(t))=K(J(t),\gamma'(t))
  \quad(0\leq t\leq t_0).
  \]
\end{theorem}
\begin{proof}
  \sketch
  By \cref{prop:dc-ch5-3-6}, the initial inner-product condition makes the
  tangential components of $J$ and $\tilde J$ equal in length; subtracting these
  parallel components reduces to the normal case. If both initial derivatives
  vanish the result is immediate. Otherwise put $v=|J|^2$ and
  $\tilde v=|\tilde J|^2$. The denominator is nonzero on $(0,a]$, and
  $v/\tilde v\to1$ at $0$. For $t_0$ with $v(t_0)\ne0$, normalize
  $U=J/\sqrt{v(t_0)}$ and $\tilde U=\tilde J/\sqrt{\tilde v(t_0)}$. The Jacobi
  equation gives
  \[
  \frac{v'(t_0)}{v(t_0)}=2I_{t_0}(U,U),\qquad
  \frac{\tilde v'(t_0)}{\tilde v(t_0)}=2I_{t_0}(\tilde U,\tilde U).
  \]
  Choose parallel orthonormal frames with first vectors tangent to the geodesics,
  and identify coefficient fields by an isometry $\phi$; then
  $\langle\phi V_1,\phi V_2\rangle=\langle V_1,V_2\rangle$ and
  $(\phi V)'=\phi(V')$. The curvature inequality gives
  $I_{t_0}(\phi U,\phi U)\leq I_{t_0}(U,U)$, while the index lemma gives
  $I_{t_0}(\tilde U,\tilde U)\leq I_{t_0}(\phi U,\phi U)$. Hence
  $(v/\tilde v)'\geq0$ and $|\tilde J|\leq|J|$. Equality at $t_0$ forces equality in
  both index comparisons on every preceding interval; the curvature inequality is
  then equality on the planes generated by the Jacobi fields, and continuity gives
  the stated interval including $0$.
\end{proof}

\begin{proposition}[Bounds on the spacing of conjugate points]
  \label{prop:dc-ch10-2-4}
  \dcref{ch10:2.4}
  If $0<L\leq K\leq H$ on $M$, then consecutive conjugate points along any geodesic
  are separated by a distance $d$ satisfying
  \[
  \frac{\pi}{\sqrt H}\leq d\leq\frac{\pi}{\sqrt L}.
  \]
\end{proposition}
\begin{proof}
  \sketch
  Compare with the sphere of constant curvature $H$. A normal Jacobi field on the
  sphere has its first zero at $\pi/\sqrt H$; Rauch comparison implies that a
  conjugate point in $M$ cannot occur earlier. Comparison with the sphere of
  curvature $L$ shows that a gap larger than $\pi/\sqrt L$ would force the sphere's
  first conjugate point to occur later, a contradiction.
\end{proof}

\begin{proposition}[Length decrease under radial curvature comparison]
  \label{prop:dc-ch10-2-5}
  \dcref{ch10:2.5}
  Let $M^n,\widetilde M^n$ satisfy $\widetilde K\geq K$ for all sectional planes.
  Fix $p\in M$, $\tilde p\in\widetilde M$, a linear isometry
  $i:T_pM\to T_{\tilde p}\widetilde M$, and $r>0$ such that $\exp_p$ is a
  diffeomorphism on $B_r(0)$ and $\exp_{\tilde p}$ is nonsingular on the corresponding
  ball. For a differentiable curve $c$ in $\exp_p(B_r(0))$, define
  \[
  \tilde c(s)=\exp_{\tilde p}\!\left(i\bigl(\exp_p^{-1}(c(s))\bigr)\right).
  \]
  Then $\ell(c)\geq\ell(\tilde c)$.
\end{proposition}
\begin{proof}
  \sketch
  Write $\bar c=\exp_p^{-1}c$ and consider the radial surface
  $f(t,s)=\exp_p(t\bar c(s))$. Its variational field $J_s$ along each radial geodesic
  has $J_s(0)=0$, $J_s(1)=c'(s)$, and initial derivative $\bar c'(s)$. The analogous
  surface in $\widetilde M$ has field $\tilde J_s$ with endpoint $\tilde c'(s)$ and
  initial derivative $i\bar c'(s)$. The initial lengths and tangential components
  agree, so Rauch gives $|\tilde c'(s)|\leq|c'(s)|$. Integrating over $s$ proves the
  length inequality.
\end{proof}

\begin{remark}[Ricci lower bound and volume comparison]
  \label{rem:dc-ch10-2-6}
  \dcref{ch10:2.6}
  If $M$ is complete with $\operatorname{Ric}_M\geq H$ and
  $\widetilde M(H)$ is the complete simply connected space of constant sectional
  curvature $H$, then normal balls satisfy
  \[
  \operatorname{vol} B_r(p)\leq \operatorname{vol}B_r(\tilde p).
  \]
  Equality implies $\operatorname{Ric}_M(\gamma'(t))=H$ for every normalized radial
  geodesic and $t<r$. The analogous upper-Ricci statement is false in general.
\end{remark}

\begin{remark}[Global volume-ratio monotonicity]
  \label{rem:dc-ch10-2-7}
  \uses{rem:dc-ch9-3-6, rem:dc-ch10-2-6}
  \dcref{ch10:2.7}
  With the notation of \cref{rem:dc-ch10-2-6}, set
  $V_r=\operatorname{vol}B_r(p)$ and $\widetilde V_r=\operatorname{vol}B_r(\tilde p)$.
  If $\operatorname{Ric}_M\geq H$, then for $R\geq r>0$
  \[
  \frac{V_r}{\widetilde V_r}\geq\frac{V_R}{\widetilde V_R}.
  \]
  For $R\leq\operatorname{diam}(M)$, equality holds exactly when
  $B_R(p)$ is isometric to $B_R(\tilde p)$.
\end{remark}

\section{Applications of the Index Lemma to immersions}

\begin{theorem}[Moore's nonimmersion theorem]
  \label{thm:dc-ch10-3-1}
  \uses{lem:dc-ch10-3-2, lem:dc-ch10-3-3}
  \dcref{ch10:3.1}
  Let $\overline M$ be complete and simply connected with
  $\overline K\leq b\leq0$, and let $M$ be compact with
  $K-\overline K\leq-b$. If $\dim\overline M<2\dim M$, no isometric immersion
  $f:M\to\overline M$ exists.
\end{theorem}
\begin{proof}
  \sketch
  Assume an immersion exists. Choose $\bar p\notin f(M)$ and $q\in M$ maximizing
  distance from $\bar p$, and let $\gamma:[0,\ell]\to\overline M$ be the normalized
  minimizing geodesic from $\bar p$ to $q$. It is perpendicular to $f(M)$ at $q$.
  For each unit $v\in T_qM$, vary $q$ in a curve tangent to $v$ and join the varied
  points radially to $\bar p$. The variational field $V$ satisfies $V(0)=0$, $V(\ell)=v$.
  Comparing the index form with the constant-curvature-$b$ model and using the index
  lemma gives $I_\ell(V,V)\geq I_\ell(\tilde J,\tilde J)$, where for a parallel field
  $\tilde w$ with $\tilde w(\ell)=\tilde v$ the model Jacobi field is
  \[
  \tilde J(t)=\frac{\sinh(t\sqrt{-b})}{\sinh(\ell\sqrt{-b})}\tilde w(t)\quad(b<0),
  \qquad \tilde J(t)=\frac{t}{\ell}\tilde w(t)\quad(b=0),
  \]
  and $I_\ell(\tilde J,\tilde J)=\sqrt{-b}\coth(\ell\sqrt{-b})>\sqrt{-b}$
  when $b<0$ (while it is positive for $b=0$). Since the
  varied radial geodesics are no longer than $\gamma$, the second variation yields
  \[
  0\geq\tfrac12E''(0)=I_\ell(V,V)+
  \langle S_{\gamma'(\ell)}v,v\rangle,
  \]
  hence
  \[
  \|B(v,v)\|>\sqrt{-b}\quad\text{for every unit }v. \tag{6}
  \]
  The Gauss equation gives, for orthonormal $v,w$,
  \[
  \langle B(v,v),B(w,w)\rangle-\|B(v,w)\|^2
  =K(v,w)-\overline K(v,w)\leq-b. \tag{7}
  \]
  Lemma \cref{lem:dc-ch10-3-3} says these inequalities force the codimension to be
  at least $\dim M$, contradicting $\dim\overline M<2\dim M$.
\end{proof}

\begin{lemma}[Second variation with the second fundamental form]
  \label{lem:dc-ch10-3-2}
  \dcref{ch10:3.2}
  If $E(s)$ is the energy of the radial variation in the proof of
  \cref{thm:dc-ch10-3-1}, then
  \[
  \tfrac12E''(0)=I_\ell(V,V)+
  \langle S_{\gamma'(\ell)}V(\ell),V(\ell)\rangle,
  \]
  where $S_{\gamma'(\ell)}$ is the shape operator of the immersion at $q$ in the
  normal direction $\gamma'(\ell)$.
\end{lemma}
\begin{proof}
  \sketch
  The second variation formula gives the index form plus
  $\langle\overline\nabla_{\overline V}\overline V,\overline N\rangle$ at $q$, for
  extensions $\overline V,\overline N$ of $V(\ell),\gamma'(\ell)$ with zero inner
  product. Metric compatibility converts this term to
  $\langle S_{\gamma'(\ell)}V(\ell),V(\ell)\rangle$.
\end{proof}

\begin{lemma}[Otsuki's bilinear-form lemma]
  \label{lem:dc-ch10-3-3}
  \dcref{ch10:3.3}
  Let $B:\mathbb R^n\times\mathbb R^n\to\mathbb R^k$ be symmetric and let $b\leq0$.
  If for every orthonormal pair $v,w$
  \[
  \langle B(v,v),B(w,w)\rangle-\|B(v,w)\|^2\leq-b,
  \qquad \|B(v,v)\|>\sqrt{-b},
  \]
  then $k\geq n$.
\end{lemma}
\begin{proof}
  \sketch
  If $k<n$, minimize $F(v)=\|B(v,v)\|^2$ on the unit sphere. At a minimizer $v$,
  differentiation along $v\cos t+w\sin t$ gives
  \[
  0=4\langle B(v,w),B(v,v)\rangle, \tag{8}
  \]
  and
  \[
  0\leq 8\|B(v,w)\|^2-4\|B(v,v)\|^2+
  4\langle B(w,w),B(v,v)\rangle. \tag{9}
  \]
  The map $w\mapsto B(v,w)$ on $v^\perp$ has image orthogonal to $B(v,v)$ and
  therefore a nonzero kernel vector $w_0$. Substituting $B(v,w_0)=0$ into (9) gives
  \[
  0\leq\langle B(w_0,w_0),B(v,v)\rangle-\|B(v,v)\|^2
  <\langle B(w_0,w_0),B(v,v)\rangle-(\sqrt{-b})^2
  \leq-b-(\sqrt{-b})^2=0,
  \]
  a contradiction.
\end{proof}

\begin{corollary}[Tompkins nonimmersion corollary]
  \label{cor:dc-ch10-3-4}
  \uses{rem:dc-ch10-2-7}
  \dcref{ch10:3.4}
  A compact flat $n$-manifold has no isometric immersion in $\mathbb R^{n+k}$ when
  $k<n$; in particular, the flat torus $T^n$ cannot be immersed in
  $\mathbb R^{2n-1}$.
\end{corollary}

\begin{corollary}[O'Neill nonimmersion corollary]
  \label{cor:dc-ch10-3-5}
  \dcref{ch10:3.5}
  If $\overline M$ is complete simply connected with $K_{\overline M}\leq0$, $M$ is
  compact with $\dim\overline M<2\dim M$, and $K_M\leq K_{\overline M}$, then no
  isometric immersion $M\to\overline M$ exists.
\end{corollary}

\begin{remark}[Necessity of compactness]
  \label{rem:dc-ch10-3-6}
  \uses{thm:dc-ch10-3-1}
  \dcref{ch10:3.6}
  Compactness in \cref{thm:dc-ch10-3-1} cannot be dropped: complete nonpositively
  curved surfaces can immerse in $\mathbb R^3$. The hyperbolic plane, however, has
  no isometric immersion in $\mathbb R^3$.
\end{remark}

\begin{remark}[Chern--Kuiper generalization]
  \label{rem:dc-ch10-3-7}
  \uses{thm:dc-ch10-3-1}
  \dcref{ch10:3.7}
  Let $M^n$ be compact and suppose each $T_pM$ contains an $m$-dimensional subspace
  $V$ on which all sectional curvatures are nonpositive. If $k<m$, there is no
  isometric immersion $M^n\to\mathbb R^{n+k}$.
\end{remark}

\begin{remark}[Failure without simple connectivity]
  \label{rem:dc-ch10-3-8}
  \uses{thm:dc-ch10-3-1, rem:dc-ch10-2-7}
  \dcref{ch10:3.8}
  The flat torus $T^{n+1}$ contains $T^n$ as a totally geodesic submanifold, so the
  simple-connectivity hypothesis on the ambient manifold in
  \cref{thm:dc-ch10-3-1} is necessary.
\end{remark}

\section{Focal points and an extension of Rauch's theorem}

Let $N\subset M$ be a submanifold and let $\gamma$ start orthogonally from
$p\in N$. Variations through geodesics whose initial points lie in $N$ and whose
initial velocities are normal to $N$ produce Jacobi fields satisfying the boundary
conditions below.

\begin{lemma}[Boundary conditions for normal variations]
  \label{lem:dc-ch10-4-1}
  \dcref{ch10:4.1}
  For such a variation, its Jacobi field $J$ satisfies
  \begin{enumerate}
  \item $J(0)\in T_pN$;
  \item $J'(0)+S_{\gamma'(0)}(J(0))\in(T_pN)^\perp$,
  \end{enumerate}
  where $S_{\gamma'(0)}$ is the shape operator of $N$. Conversely, every Jacobi
  field satisfying these two conditions is the variational field of a variation
  through geodesics normal to $N$ at their initial points.
\end{lemma}
\begin{proof}
  \sketch
  The first condition is $J(0)=\alpha'(0)$ for the initial curve $\alpha$ in $N$.
  Differentiating the normality of the initial velocity $A(s)$ gives, for every
  $v\in T_pN$,
  $\langle J'(0),v\rangle=-\langle S_{\gamma'(0)}J(0),v\rangle$, proving the
  second condition. Conversely, choose $\alpha$ in $N$ with $\alpha'(0)=J(0)$,
  extend $\gamma'(0)$ to a field $W$ along $\alpha$ with prescribed covariant
  derivative $J'(0)$, project $W$ to the normal bundle, and set
  $f(s,t)=\exp_{\alpha(s)}(tV(s))$. The boundary condition ensures that its
  variational field agrees with $J$.
\end{proof}

\begin{definition}[Focal point]
  \label{def:dc-ch10-4-2}
  \dcref{ch10:4.2}
  A point $q$ is focal to $N\subset M$ if there are $p\in N$ and a geodesic
  $\gamma:[0,\ell]\to M$ with $\gamma(0)=p$, $\gamma'(0)\perp T_pN$, and
  $\gamma(\ell)=q$, together with a nonzero Jacobi field satisfying the two
  conditions of \cref{lem:dc-ch10-4-1} and $J(\ell)=0$.
\end{definition}

\begin{example}[Focal points of an equator]
  \label{ex:dc-ch10-4-3}
  \dcref{ch10:4.3}
  The north and south poles are focal points of the equator $S^{n-1}\subset S^n$.
  More generally, for the normal bundle $T(N)^\perp\to N$, the normal exponential
  map is the restriction
  \[
  \exp^\perp:T(N)^\perp\to M,
  \]
  and $\dim T(N)^\perp=\dim M$.
\end{example}

\begin{proposition}[Focal points and the normal exponential map]
  \label{prop:dc-ch10-4-4}
  \uses{lem:dc-ch10-4-1}
  \dcref{ch10:4.4}
  A point $q$ is focal to $N$ if and only if it is a critical value of
  $\exp^\perp:T(N)^\perp\to M$.
\end{proposition}
\begin{proof}
  \sketch
  A normal geodesic variation with a focal Jacobi field gives a curve
  $w(s)=(\alpha(s),\ell A(s))$ in $T(N)^\perp$ whose image under $\exp^\perp$ has
  zero derivative at $s=0$; hence its endpoint is a critical value. Conversely,
  represent a critical tangent vector to $\exp^\perp$ by a curve
  $w(s)=(\sigma(s),\ell V(s))$, use
  $f(s,t)=\exp_{\sigma(s)}(tV(s))$, and apply
  \cref{lem:dc-ch10-4-1}. The resulting Jacobi field is nonzero, satisfies the
  boundary conditions, and vanishes at the critical endpoint.
\end{proof}

\begin{definition}[Focal set]
  \label{def:dc-ch10-4-5}
  \dcref{ch10:4.5}
  The focal set $F(N)\subset M$ is the set of focal points of $N$.
\end{definition}

\begin{example}[Focal surfaces in Euclidean space]
  \label{ex:dc-ch10-4-6}
  \uses{prop:dc-ch10-4-4}
  \dcref{ch10:4.6}
  For a regular surface $N^2\subset\mathbb R^3$, the focal set is obtained by
  moving along each unit normal a distance equal to the reciprocal of a principal
  curvature. In local coordinates with first fundamental matrix $(g_{ij})$ and
  second fundamental matrix $(h_{ij})$, the normal exponential map is
  $\mathbf x+tn$ and is singular exactly when $(g_{ij}-th_{ij})$ is singular. Thus
  $\mathbf x+tn$ is focal precisely when $1/t$ is a principal curvature. For a
  sphere the focal set is its center; for a circular cylinder it is its axis.
\end{example}

\begin{definition}[Focal-point-free geodesic]
  \label{def:dc-ch10-4-7}
  \dcref{ch10:4.7}
  A geodesic $\gamma:[0,a]\to M$ is focal-point free on $(0,a]$ if some ball
  $B_\varepsilon(0)\subset\gamma'(0)^\perp$ has the property that $\gamma$ has no
  focal point relative to $\Sigma_\varepsilon=\exp_{\gamma(0)}B_\varepsilon(0)$.
  Since $\Sigma_\varepsilon$ is geodesic at $\gamma(0)$, a Jacobi field with
  $J(0)\ne0$ and $J'(0)=0$ has $S_{\gamma'(0)}J(0)=0$.
\end{definition}

\begin{lemma}[Index comparison for focal-point-free geodesics]
  \label{lem:dc-ch10-4-8}
  \dcref{ch10:4.8}
  If $\gamma:[0,a]\to M^n$ is focal-point free on $(0,a]$, $J$ is a normal Jacobi
  field with $J'(0)=0$, and $V$ is piecewise differentiable with
  $J(t_0)=V(t_0)$, then
  \[
  I_{t_0}(J,J)\leq I_{t_0}(V,V),
  \]
  with equality exactly when $V=J$ on $[0,t_0]$.
\end{lemma}
\begin{proof}
  \sketch
  Jacobi fields normal to $\gamma$ with derivative zero at $0$ form an
  $(n-1)$-dimensional space. Focal-point freeness makes their values a basis of
  $\gamma'(t)^\perp$ for every $t>0$. Expanding $V$ in this basis and repeating the
  integration-by-parts argument of \cref{lem:dc-ch10-2-2} yields the inequality and
  the same equality criterion.
\end{proof}

\begin{theorem}[Rauch comparison with focal points]
  \label{thm:dc-ch10-4-9}
  \uses{lem:dc-ch10-4-8, thm:dc-ch10-2-3}
  \dcref{ch10:4.9}
  Let $\gamma:[0,a]\to M^n$ and $\tilde\gamma:[0,a]\to\widetilde M^{n+k}$ have
  equal speeds, and let $J,\tilde J$ be Jacobi fields satisfying
  \[
  J'(0)=\tilde J'(0)=0,\quad
  \langle J(0),\gamma'(0)\rangle=\langle\tilde J(0),\tilde\gamma'(0)\rangle,
  \quad |J(0)|=|\tilde J(0)|.
  \]
  Assume $\tilde\gamma$ is focal-point free on $(0,a]$ and
  \[
  \widetilde K(\tilde x,\tilde\gamma'(t))\geq K(x,\gamma'(t))
  \]
  for all $t$ and all tangent vectors $x,\tilde x$. Then
  $|\tilde J(t)|\leq|J(t)|$. Equality at $t_0>0$ implies
  \[
  \widetilde K(\tilde J(t),\tilde\gamma'(t))=K(J(t),\gamma'(t))
  \quad(0\leq t\leq t_0).
  \]
\end{theorem}
\begin{proof}
  \sketch
  Repeat the quotient-of-squared-lengths argument in
  \cref{thm:dc-ch10-2-3}, now with initial derivatives zero. Focal-point freeness
  ensures that $|J|^2/|\tilde J|^2$ is defined on $(0,a]$, and the focal index lemma
  reverses the comparison step to give
  $I_{t_0}(\phi(U),\phi(U))\geq I_{t_0}(\tilde U,\tilde U)$. The resulting quotient is
  nondecreasing, proving the norm inequality; equality forces equality in the index
  and curvature comparisons on every preceding interval, yielding the curvature
  equality by continuity.
\end{proof}
